\documentclass{article}

% if you need to pass options to natbib, use, e.g.:
%     \PassOptionsToPackage{numbers, compress}{natbib}
% before loading neurips_2025

% The authors should use one of these tracks.
% Before accepting by the NeurIPS conference, select one of the options below.
% 0. "default" for submission
 \usepackage{neurips_2025}
% the "default" option is equal to the "main" option, which is used for the Main Track with double-blind reviewing.
% 1. "main" option is used for the Main Track
%  \usepackage[main]{neurips_2025}
% 2. "position" option is used for the Position Paper Track
%  \usepackage[position]{neurips_2025}
% 3. "dandb" option is used for the Datasets & Benchmarks Track
 % \usepackage[dandb]{neurips_2025}
% 4. "creativeai" option is used for the Creative AI Track
%  \usepackage[creativeai]{neurips_2025}
% 5. "sglblindworkshop" option is used for the Workshop with single-blind reviewing
 % \usepackage[sglblindworkshop]{neurips_2025}
% 6. "dblblindworkshop" option is used for the Workshop with double-blind reviewing
%  \usepackage[dblblindworkshop]{neurips_2025}

% After being accepted, the authors should add "final" behind the track to compile a camera-ready version.
% 1. Main Track
 % \usepackage[main, final]{neurips_2025}
% 2. Position Paper Track
%  \usepackage[position, final]{neurips_2025}
% 3. Datasets & Benchmarks Track
 % \usepackage[dandb, final]{neurips_2025}
% 4. Creative AI Track
%  \usepackage[creativeai, final]{neurips_2025}
% 5. Workshop with single-blind reviewing
%  \usepackage[sglblindworkshop, final]{neurips_2025}
% 6. Workshop with double-blind reviewing
%  \usepackage[dblblindworkshop, final]{neurips_2025}
% Note. For the workshop paper template, both \title{} and \workshoptitle{} are required, with the former indicating the paper title shown in the title and the latter indicating the workshop title displayed in the footnote.
% For workshops (5., 6.), the authors should add the name of the workshop, "\workshoptitle" command is used to set the workshop title.
% \workshoptitle{WORKSHOP TITLE}

% "preprint" option is used for arXiv or other preprint submissions
 % \usepackage[preprint]{neurips_2025}

% to avoid loading the natbib package, add option nonatbib:
%    \usepackage[nonatbib]{neurips_2025}

\usepackage[utf8]{inputenc} % allow utf-8 input
\usepackage[T1]{fontenc}    % use 8-bit T1 fonts
\usepackage{hyperref}       % hyperlinks
\usepackage{url}            % simple URL typesetting
\usepackage{booktabs}       % professional-quality tables
\usepackage{amsfonts}       % blackboard math symbols
\usepackage{nicefrac}       % compact symbols for 1/2, etc.
\usepackage{microtype}      % microtypography
\usepackage{xcolor}         % colors
\usepackage{amsmath}
\usepackage{graphicx}
\usepackage{algorithm}
\usepackage{algorithmic}
% Note. For the workshop paper template, both \title{} and \workshoptitle{} are required, with the former indicating the paper title shown in the title and the latter indicating the workshop title displayed in the footnote. 
\title{STED and Consistency Scoring: A Framework for Evaluating LLM Structured Output Reliability}


% The \author macro works with any number of authors. There are two commands
% used to separate the names and addresses of multiple authors: \And and \AND.
%
% Using \And between authors leaves it to LaTeX to determine where to break the
% lines. Using \AND forces a line break at that point. So, if LaTeX puts 3 of 4
% authors names on the first line, and the last on the second line, try using
% \AND instead of \And before the third author name.


\author{%
  David S.~Hippocampus\thanks{Use footnote for providing further information
    about author (webpage, alternative address)---\emph{not} for acknowledging
    funding agencies.} \\
  Department of Computer Science\\
  Cranberry-Lemon University\\
  Pittsburgh, PA 15213 \\
  \texttt{hippo@cs.cranberry-lemon.edu} \\
  % examples of more authors
  % \And
  % Coauthor \\
  % Affiliation \\
  % Address \\
  % \texttt{email} \\
  % \AND
  % Coauthor \\
  % Affiliation \\
  % Address \\
  % \texttt{email} \\
  % \And
  % Coauthor \\
  % Affiliation \\
  % Address \\
  % \texttt{email} \\
  % \And
  % Coauthor \\
  % Affiliation \\
  % Address \\
  % \texttt{email} \\
}


\begin{document}


\maketitle


\begin{abstract}
Large Language Models (LLMs) are increasingly deployed for structured data generation, yet output consistency remains critical for production applications. We introduce a comprehensive framework for evaluating and improving consistency in LLM-generated structured outputs. Our approach combines: (1) STED (Semantic Tree Edit Distance), a novel similarity metric balancing semantic flexibility with structural strictness when comparing JSON outputs, and (2) a consistency scoring framework aggregating multiple STED measurements across repeated generations to quantify reliability. Through systematic experiments on synthetic datasets with controlled schema, expression, and semantic variations, we demonstrate STED achieves superior performance ($0.86-0.90$ similarity for semantic equivalents, $0.0$ for structural breaks) compared to existing metrics including TED, BERTScore, and DeepDiff. Applying our framework to benchmark six LLMs reveals significant variations: Claude-3.7-Sonnet demonstrates exceptional consistency, maintaining near-perfect structural reliability even at high temperatures ($T=0.9$), while models like Claude-3-Haiku and Nova-Pro exhibit substantial degradation requiring careful tuning. Our framework enables practical applications including targeted model selection for structured tasks, iterative prompt refinement for reproducible results, and diagnostic analysis to identify inconsistency root causes. This work provides theoretical foundations and practical tools for ensuring reliable structured output generation in LLM-based production systems.
\end{abstract}


\section{Introduction}

Large Language Models (LLMs) have become integral to production systems requiring structured data generation, particularly in JSON format for APIs, data extraction pipelines, and automated workflows. However, evaluating the consistency and quality of LLM-generated structured outputs presents unique challenges that existing evaluation methods fail to adequately address.

A fundamental issue in evaluating structured outputs is the mismatch between evaluation methods and structured semantics. Consider two structured objects with identical content but different key ordering---while the structured specification treats these as equivalent, popular evaluation methods like BERTScore \cite{Zhang2020} assign them significantly different similarity scores due to position-sensitive embeddings, leading to false negatives in production systems.

The challenge extends beyond simple key reordering. LLMs frequently generate functionally equivalent structures with various forms of semantic equivalence: naming convention differences (\texttt{"user\_name"} versus \texttt{"userName"}), array reorderings where sequence is not semantically significant, structural reorganizations preserving information content, and type representation variations (\texttt{"true"} versus \texttt{true}, \texttt{"123"} versus \texttt{123}). Each represents valid structured data conveying identical information, yet existing metrics fail to recognize their equivalence.

Current evaluation approaches exhibit systematic limitations when applied to structured outputs. BERTScore \cite{Zhang2020}, while effective for natural language evaluation, suffers from order sensitivity that violates JSON's order-agnostic semantics for object properties, systematically underestimating similarity for reordered but identical structures. Tree Edit Distance (TED) methods focus purely on structural differences without considering semantic equivalence, while exact matching approaches are overly restrictive for practical applications.

\textbf{DeepDiff} \cite{seperman} and similar structural comparison tools can ignore ordering through configuration but lack semantic understanding. They treat \texttt{"email"} and \texttt{"email\_address"} as completely different keys, missing obvious semantic relationships that humans would recognize.

These limitations have significant practical consequences. In production systems processing thousands of API responses, false negatives from order sensitivity can trigger unnecessary alerts, increase operational overhead, and mask genuine inconsistencies.

While traditional edit distances like Tree Edit Distance (TED) have been applied to structured data comparison, they fail to handle the semantic equivalences common in LLM outputs. Recent advances in learned edit distances \cite{paassen2018tree} improve classification accuracy but focus on discriminative tasks rather than consistency evaluation. We propose STED (Semantic Tree Edit Distance), which adapts edit distance specifically for LLM output evaluation, balancing semantic equivalence with structural validity.

Unlike embedding-based approaches that learn vectorial representations \cite{paassen2018tree}, STED directly incorporates JSON semantics through three key innovations: (1) \textbf{Semantic-Enhanced Tree Edit Distance} that recognizes semantically equivalent keys and values while preserving structural constraints, (2) \textbf{Order-Invariant Matching} using Hungarian algorithm for optimal element pairing, and (3) \textbf{Multi-Level Similarity} integrating structural, key, value, and type similarities with configurable weights.

Our experiments demonstrate STED's superior discrimination: while BERTScore and DeepDiff score $>0.95$ across all variations---failing to distinguish schema violations from benign reorderings---STED correctly identifies critical differences ($0.23$ for schema violations vs $>0.95$ for order variations), achieving $4\times$ better discrimination.

As LLMs increasingly power production systems, STED provides a theoretically grounded solution for evaluating structured output consistency, enabling reliable deployment through accurate distinction between critical errors and acceptable variations.
\vspace{-3mm}
\section{Related Work}
\vspace{-3mm}
\subsection{Edit Distance and Structured Data Comparison}
The tree edit distance problem, formalized by Tai \cite{Tai1979}, provides the foundation for comparing hierarchical structures. Zhang and Shasha \cite{Zhang1989} proposed the first polynomial-time algorithm, later improved by RTED \cite{Pawlik2011} and extended with move operations \cite{Chawathe1996}. However, these classical approaches rely on exact node matching without semantic understanding. Recent learning-based approaches \cite{paassen2018tree, paassen2022revisiting} explore adaptive embeddings, while graph edit distance methods \cite{jain2024graph, piao2023computing} handle complex structures, inspiring our semantic enhancements.
JSON-specific tools like RFC 6902 \cite{RFC6902} and DeepDiff \cite{seperman} provide structural comparison with optional order-insensitive matching but lack semantic awareness. JSON Schema validation \cite{Wright2022} focuses on conformance rather than similarity, while query languages like JMESPath \cite{JMESPath2014} and JSONPath \cite{Gessner2007} target extraction, not comparison.
\subsection{Neural Approaches to Similarity}
BERTScore \cite{Zhang2020} revolutionized text evaluation using contextualized embeddings from BERT \cite{Devlin2019}, but its position sensitivity makes it unsuitable for JSON where key ordering is irrelevant. Advances like Sentence-BERT \cite{Reimers2019}, SimCSE \cite{Gao2021}, and Universal Sentence Encoder \cite{Cer2018} improve embedding quality but process JSON as flat strings, losing structural information. Code similarity work including CodeBERT \cite{Feng2020} and GraphCodeBERT \cite{Guo2021} demonstrates the importance of combining structural and semantic features. Graph neural networks offer another perspective through GraphSAGE \cite{Hamilton2017}, GAT \cite{Velickovic2018}, Tree-LSTM \cite{Tai2015}, and GraphFormers \cite{Yang2021}. While powerful, these require training data and are computationally expensive compared to our approach.
\subsection{LLM Structured Output Generation and Consistency}
Recent work addresses structured output challenges in LLMs. Bubeck et al. \cite{Bubeck2023} analyze GPT-4's structured generation capabilities, while HELM \cite{Liang2022} provides holistic evaluation focusing on task performance. SLOT \cite{Wang2024slot} and StructuredRAG \cite{shorten2024structuredrag} address output structuring, with StructEval \cite{yang2025structeval} benchmarking generation capabilities. Format control studies \cite{wang-etal-2025-verifiable, geng2025jsonschema} examine verification and schema conformance.
LLM consistency research spans multiple dimensions. Elazar et al. \cite{elazar2021measuring} and Raj et al. \cite{Raj2023} examine consistency across paraphrases, while recent studies investigate stability \cite{atil2024llm}, non-determinism \cite{baldwin2024nondeterminism}, and automated analysis \cite{patwardhan2025automated, wu2025estimating}. Temperature effects \cite{renze2024effect, renze-2024-effect} and prompt stability \cite{chen2025prompt} are also explored, but these focus on input variations rather than output consistency for identical inputs.
\subsection{Semantic Matching and Optimal Assignment}
Semantic similarity in structured data draws from ontology matching \cite{Euzenat2013} and schema matching \cite{Rahm2001}, with approaches like SimFlood \cite{Melnik2002} and COMA \cite{Do2002} combining multiple strategies. Neural methods including DeepMatcher \cite{Mudgal2018} and SMAT \cite{Zhang2021smat} achieve high accuracy but target one-time alignment rather than continuous evaluation. Distance functions for structured data \cite{ontanon2020overview, kriegel2005similarity, yang2005similarity} provide theoretical foundations.
The Hungarian algorithm \cite{Kuhn1955, Munkres1957} enables polynomial-time optimal assignment, extended recently for large-scale approximations \cite{Schwiegelshohn2022}. TreeKernel \cite{Collins2001} and subsequent work \cite{Moschitti2006} apply assignment to tree matching, inspiring our combination with semantic similarity for arrays and keys in JSON.

Despite extensive related work, existing methods fail to address structured LLM-generated outputs consistency evaluation challenges:
(1) Methods focus on either structure or semantics but not both;
(2) Tools are either order-agnostic or overly sensitive;
(3) No method handles LLM-specific variation patterns.
STED bridges these gaps through unified structural-semantic analysis with appropriate order handling and granular insights.
\vspace{-3mm}
\section{Methodology}
\vspace{-3mm}
We present STED (Semantic Tree Edit Distance), a novel framework for evaluating consistency in LLM-generated structured outputs. While traditional metrics treat structured formats like JSON, XML, and HTML as flat text, we recognize their hierarchical nature by transforming them into tree representations. This allows us to reframe the structured output consistency problem as tree consistency evaluation. Our approach extends classical tree edit distance algorithms with semantic similarity capabilities to compute pairwise distances, then aggregates these distances across multiple LLM generations into a normalized consistency score. This two-stage process addresses the fundamental challenge of quantifying output reliability when LLMs produce functionally equivalent but syntactically different structures.

\subsection{Problem Formulation}

Let $\mathcal{O} = \{o_1, o_2, ..., o_n\}$ denote structured outputs generated by an LLM for identical inputs. While our framework applies to any hierarchical format (JSON, XML, HTML), we use JSON for clarity. Each output $o_i$ is represented as a tree $T_i = (V_i, E_i)$, where $V_i$ contains nodes (keys, values, structural elements) and $E_i$ encodes parent-child relationships.

The consistency evaluation problem involves two stages:

\textbf{Pairwise Similarity:} Given trees $T_1$ and $T_2$, compute a similarity score $s(T_1, T_2) \in [0, 1]$ that captures semantic and structural equivalence, where $s = 1$ indicates perfect consistency and $s = 0$ indicates fundamental incompatibility.

\textbf{Consistency Score:} Given $n$ outputs, aggregate pairwise similarities into a global consistency score $C(\mathcal{O}) \in [0, 1]$ that quantifies the LLM's reliability:

\begin{equation}
C(\mathcal{O}) = \frac{2}{n(n-1)} \sum_{i=1}^{n-1} \sum_{j=i+1}^{n} s(T_i, T_j)
\end{equation}

The core challenge is distinguishing benign variations from critical differences that affect functionality, while providing an interpretable consistency metric for production deployment.

\subsection{Tree Representation of Structured Data}

We transform structured documents into trees $T = (V, E)$ where each node $v \in V$ contains: \textbf{Type}: $\{\text{object}, \text{array}, \text{string}, \text{number}, \text{boolean}, \text{null}\}$; \textbf{Label}: Key name for object properties; \textbf{Value}: Raw value for primitives, child list for arrays; \textbf{Path}: Hierarchical path for unique identification.

The transformation rules are: (1) \textbf{Objects} become internal nodes with labeled edges to child values, (2) \textbf{Arrays} remain as nodes containing element lists, recursively transformed if complex, and (3) \textbf{Primitives} map to leaf nodes with actual values. This preserves array structure while enabling order-invariant matching and maintaining the semantic distinction between arrays and objects.

\subsection{Semantic Tree Edit Distance}

Our STED algorithm extends classical tree edit distance with semantic awareness through three operations: \textbf{Insert} (cost $\gamma_{ins}(v)$), \textbf{Delete} (cost $\gamma_{del}(v)$), and \textbf{Update} (cost $\gamma_{upd}(v_1, v_2)$). The semantic update cost combines multiple dimensions:
\vspace{-2mm}
\begin{equation}
\gamma_{upd}(v_1, v_2) = w_s \cdot \gamma_{struct}(v_1, v_2) + w_c \cdot \gamma_{content}(v_1, v_2)
\end{equation}

where $\gamma_{struct}$ measures structural similarity using embedding-based comparison, and $\gamma_{content}$ evaluates value similarity with type-aware costs. For arrays and objects, we employ Hungarian algorithm for optimal child matching, ensuring order-invariant comparison.

\subsection{Semantic Similarity Computation}

Field names are normalized (e.g., \texttt{"userName"} $\rightarrow$ \texttt{"user name"}) to capture semantic relationships. We compute similarity using embeddings with cosine similarity for texts $<300$ characters; longer texts are chunked recursively with 50-character overlaps. We used Amazon Titan Text Embeddings v2 in our implementation, though the choice of embedding model has minimal impact on STED's performance since we primarily match short field names and values where most modern embedding models achieve similar semantic understanding. The framework remains model-agnostic and can use any embedding model (e.g., all-MiniLM-L6-v2, Sentence-BERT, OpenAI embeddings). Type-aware comparison ensures appropriate metrics: semantic for strings, exact for numbers, order-invariant for arrays. This recognizes functional equivalence (e.g., \{\texttt{"user\_name"}: \texttt{"John"}\} $\equiv$ \{\texttt{"userName"}: \texttt{"John"}\}) while detecting type violations that break compatibility.

\subsection{Optimal Subtree Matching}

Traditional tree edit distance algorithms process nodes sequentially, potentially missing globally optimal alignments. We address this through Hungarian algorithm-based matching at each tree level.

For trees $T_1$ and $T_2$ with children sets $C_1$ and $C_2$, we formulate optimal matching as an assignment problem:
\vspace{-2mm}
\begin{equation}
\text{OptimalCost}(T_1, T_2) = \min_{\pi} \left[ \sum_{(i,j) \in \pi} d(c_1^i, c_2^j) + \sum_{i \notin \pi} \gamma_{del}(c_1^i) + \sum_{j \notin \pi} \gamma_{ins}(c_2^j) \right]
\end{equation}

where $\pi$ represents the matching between children, $d(\cdot, \cdot)$ is the recursive distance, and unmatched nodes incur insertion/deletion costs.

The algorithm constructs a cost matrix $M \in \mathbb{R}^{\max(|C_1|, |C_2|) \times \max(|C_1|, |C_2|)}$ where $M_{ij}$ represents the cost of matching child $i$ from $T_1$ with child $j$ from $T_2$. The Hungarian algorithm finds the minimum-cost assignment in $O(n^3)$ time, ensuring globally optimal child alignment rather than greedy local decisions.

\subsection{Similarity Score Normalization}

The STED algorithm normalizes the computed distance at each tree level to produce an interpretable similarity score in [0,1]:
\vspace{-2mm}
\begin{equation}
\text{STED}(T_1, T_2) = 1 - \min\left(1, \frac{d_{\text{matched}} + \lambda \cdot \Delta_{\text{unmatched}}}{\max(|C_1|, |C_2|)}\right)
\end{equation}

where $d_{\text{matched}}$ is the total cost from Hungarian assignment, $\Delta_{\text{unmatched}} = ||C_1| - |C_2||$ penalizes size differences with weight $\lambda = 0.1$. This per-level normalization ensures consistent scoring regardless of tree depth.

\subsection{Consistency Score Calculation}

Let $\{s_i\}_{i=1}^n$ be the similarity values between responses across different runs, with empirical standard deviation $\sigma = \operatorname{std}(s_1, \dots, s_n)$. To normalize for scale, we compute the maximum possible standard deviation for $n$ values in $[0,1]$:

\[
\sigma_{\max} = \operatorname{std}\big(\underbrace{0, \dots, 0}_{\lfloor n/2 \rfloor}, \underbrace{1, \dots, 1}_{\lceil n/2 \rceil}\big).
\]

We then define the normalized deviation

\[
\hat{\sigma} = 
\begin{cases}
\dfrac{\sigma}{\sigma_{\max}}, & \sigma_{\max} > 0, \\
0, & \text{otherwise}.
\end{cases}
\]

Finally, the consistency score is given by
\vspace{-2mm}
\[
\text{ConsistencyScore}(s_1,\dots,s_n) = \left( \frac{1}{1 + 2\hat{\sigma}} \right)^{\alpha},
\]

where $\alpha = 20$ is a steepness factor that amplifies the typically small deviations observed in model outputs, providing better discrimination in the common low-deviation range while saturating for rare large deviations. If $n \leq 1$, we set $\text{ConsistencyScore} = 1$.

\subsection{Computational Complexity}

STED requires $O(N \times B^3)$ time where $N$ is total node count and $B$ is maximum branching factor, with the Hungarian algorithm contributing the cubic factor. Space complexity is $O(B^2 + D)$ for cost matrices and recursion depth $D$. Large arrays/objects with high branching factor $B$ dominate cost, while tree depth affects only space linearly, making STED tractable for typical structured outputs but potentially expensive for highly-branched structures.

\section{Experiments and Results}
\vspace{-2mm}
Our experimental evaluation comprises two complementary studies. First, we validate our method's effectiveness using synthetic datasets with controlled variations to verify its ability to accurately quantify similarity degradation. Second, we deploy our framework to benchmark the structured output consistency of six LLMs available through Amazon Bedrock: Claude 3.7 Sonnet, Claude 3.5 Sonnet V2, Claude 3.5 Haiku, Claude 3 Haiku, Llama-3.3-70B, and Nova Pro. This model selection represents diverse architectures and optimization strategies—from efficiency-focused (Haiku variants) to performance-focused (Sonnet variants). Our goal is validating the evaluation framework's effectiveness rather than exhaustive model coverage. The chosen models sufficiently demonstrate our method's ability to reveal consistency patterns (temperature sensitivity, structural-semantic gaps, degradation profiles) that represent fundamental behaviors in autoregressive language models. The framework itself is model-agnostic and can evaluate any LLM with structured output capabilities.

\subsection{Method Effectiveness Verification}

\subsubsection{Variation Taxonomy}
LLM-generated structured data exhibits three distinct variation types: \textbf{Schema Variation}: Structural modifications including field name changes, structure flattening, and hierarchy alterations that affect parseability; \textbf{Expression Variation}: Lexical modifications preserving semantic meaning through synonym substitution, paraphrasing, and abbreviation usage; \textbf{Semantic Variation}: Fundamental content changes that alter data meaning, potentially causing incorrect interpretations.

\subsubsection{Dataset Construction}

\textbf{Our evaluation leverages 2,400 synthetic test cases} systematically generated from a diverse set of 80 base samples. We begin with 80 ShareGPT-formatted JSON samples from Quiz Generation and Structured Output datasets (75 valid after parsing errors), which serve as seeds for controlled variation generation. These base samples ensure structural diversity: depths range from 2--7 levels (mean 4.0±1.0, mode 4), field counts span 4--228 (mean 42.3±37.8), and include realistic type distributions (68.2\% strings, 18.4\% integers, 13.4\% complex types). A comprehensive distribution analysis is provided in Appendix~\ref{sec:dataset-distribution}.

\textbf{Gradual Variations (2,250 samples):} We create variants with controlled modification ratios from 0.1 to 1.0 (10 levels) for: (1) \textit{Field Name Variants} - semantically equivalent keys (e.g., "user\_name" → "userName"), (2) \textit{Expression Variants} - paraphrased values preserving meaning, and (3) \textit{Semantic Variants} - content changes affecting functionality. Each category yields 750 samples (75 × 10 levels).

\textbf{Structural Variations (150 samples):} We generate single variants for: (1) \textit{Flattened Structure} - nested-to-flat transformations potentially causing field collisions, and (2) \textit{Nested Changes} - modified hierarchies breaking API compatibility. These binary changes yield 75 samples each.

\subsubsection{Baseline Methods}
We compare STED against three established approaches: \textbf{TED}: Tree Edit Distance using Zhang-Shasha algorithm with default cost functions; \textbf{BERTScore}: Adapted for JSON by serializing structures with sorted keys, computing F1 scores from BERT embedding alignments; \textbf{DeepDiff}: Rule-based structural comparison using Deep Distance metric, converted to similarity as $1 - \text{DeepDistance}$.

\subsection{LLM Consistency Benchmarking}

We design a comprehensive framework to evaluate LLM consistency in structured output generation. For each dataset sample, we augment prompts with extracted ground truth schemas. Our protocol executes each of 75 samples 10 times per temperature (0.1-0.9), yielding \textbf{750 outputs per temperature setting} and 6,750 total outputs per model for robust consistency analysis.

We employ three evaluation modes: (1) structural consistency measuring format adherence, (2) semantic consistency evaluating content preservation, and (3) hybrid approach with equal weighting ($\alpha = 0.5$). This multi-faceted evaluation reveals how temperature variations affect different aspects of structured output generation.

\begin{figure}[t]
\centering
\includegraphics[width=1.0\textwidth]{figures/schema_variation_analysis.png}
\caption{Similarity scores across schema variation types for four consistency evaluation methods. Field name changes (0.1-1.0 ratios) show gradual similarity degradation, with BERTScore maintaining highest scores. For structural variations, STED achieves zero similarity on nested changes but moderate scores on flat structures. All methods use similarity scores ranging from 0 to 1.}
\label{fig:schema_variation_analysis}
\end{figure}

\begin{figure}[t]
\centering
\includegraphics[width=1.0\textwidth]{figures/similarity_progression_combined.png}
\caption{Similarity score progression under semantic variations (left) and expression variations (right). Ideal behavior: sensitivity to semantic changes, robustness to expression changes. STED achieves optimal balance with controlled degradation for semantic variations while maintaining high scores for expression variations, aligning with human consistency perception. (TED remains constant at 1.0 as it measures only structural similarity, not content.)}
\label{fig:similarity_progression_combined}
\end{figure}

\subsection{Method Effectiveness Verification}

\subsubsection{Schema Variation Analysis}
Figure~\ref{fig:schema_variation_analysis} reveals how methods handle structural variations in LLM outputs. For field name variations using semantic equivalents, STED maintains stable similarity across all variation ratios, correctly recognizing functional equivalence. TED degrades significantly due to lack of semantic understanding, while BERTScore's excessive permissiveness may mask important structural differences. For structural modifications (flattening and nesting changes), STED correctly assigns zero similarity, recognizing these as breaking changes for downstream systems. Other methods problematically tolerate these changes with non-zero scores that could allow incompatible outputs to pass consistency checks.

\subsubsection{Content Variation Analysis}
Figure~\ref{fig:similarity_progression_combined} evaluates content variation handling, critical for applications like content moderation where semantic accuracy impacts outcomes. For semantic variations, STED (0.954±0.039) and BERTScore (0.958±0.025) show statistically equivalent calibrated sensitivity (p=0.600), while DeepDiff over-reacts (0.799±0.165) and TED remains insensitive (1.0, structure-only). For expression variations, STED achieves superior robustness (0.981±0.017, p<0.001), correctly preserving high similarity for paraphrases, unlike DeepDiff which incorrectly penalizes valid rewording (0.805±0.164).\footnote{Detailed statistical analysis in Supplementary Section~\ref{sec:statistical-analysis}.} STED's differential response—5.2\% degradation for semantic changes versus 1.9\% for expression changes—enables threshold-based detection of genuine semantic drift while accepting natural language variation, providing human-aligned consistency assessment for practical applications.

\vspace{-3mm}
\subsection{Model Consistency Analysis}
\vspace{-3mm}
Figure~\ref{fig:consistency_score_by_consistency_type} reveals critical insights into consistency patterns across six language models, providing guidance for model selection and parameter tuning.

\subsubsection{Temperature Response Patterns}
All models exhibit inverse temperature-consistency relationships, validating our evaluation framework. Claude-3.7-Sonnet demonstrates superior stability with only 19\% degradation from T=0.1 to T=0.9, maintaining 0.658±0.240 mean consistency. Claude-3.5-Haiku shows highest sensitivity (46\% decline), while Claude-3-5-Sonnet-V2 plateaus between T=0.2-0.6, requiring higher temperatures for diversity.

\subsubsection{Structural Preservation Under Diversity}
Claude-3.7-Sonnet achieves near-perfect structural consistency (0.999±0.037) across all temperatures, demonstrating exceptional format preservation even at T=0.9. This 0.5-point structural-semantic gap indicates sophisticated behavior: varying content while preserving format—ideal for template-based generation. Claude-3-5-Sonnet-V2 similarly maintains 0.946±0.227 structural consistency while semantic scores vary more widely (0.459±0.318), confirming models can reliably produce diverse content within consistent formats.

\subsubsection{Deployment Considerations}
The evaluation reveals key selection criteria. Structural consistency ranges from near-perfect (Claude-3.7-Sonnet) to moderate (Nova-Pro: 0.728±0.430). Temperature responses vary from smooth degradation to erratic patterns (Nova-Pro peaks at T=0.5). For production: T=0.1-0.3 maximizes consistency, T=0.4-0.6 balances diversity with acceptable degradation, $T\geq0.7$ causes significant losses (>30\% semantic drop). The structural-semantic gap guides task selection: larger gaps suit template-based generation, smaller gaps indicate uniform variation for free-form tasks.

\begin{figure}[t]
\centering
\includegraphics[width=1.0\textwidth]{figures/consistency_score_by_consistency_type_with_errors.png}
\caption{Model consistency across temperature settings. Higher scores indicate better consistency between outputs for the same prompts. Error bars show standard deviation across 75 test cases per temperature. \textit{Detailed numerical results in Supplementary Tables \ref{tab:overall-consistency-temp}--\ref{tab:consistency-summary}.}}
\label{fig:consistency_score_by_consistency_type}
\end{figure}
\vspace{-3mm}
\section{Conclusion}
We present STED and consistency score, a novel metric for evaluating consistency in LLM-generated structured outputs that balances semantic flexibility with structural strictness. Through controlled experiments on synthetic datasets, we demonstrate STED's superiority: it correctly assigns zero similarity to structure-breaking modifications while maintaining robustness to semantically equivalent variations (0.86-0.90), outperforming TED, BERTScore, and DeepDiff.
Our benchmarking of six state-of-the-art LLMs reveals critical insights for production deployment. The proposed consistency score enables three key applications: (1) model selection for structured output tasks, providing targeted evaluation beyond general-purpose benchmarks; (2) prompt refinement through iterative optimization, enabling developers to craft prompts yielding reproducible outputs; and (3) diagnostic analysis to identify factors contributing to inconsistency. These capabilities make our framework a practical tool for improving the reliability of LLM-based systems in production environments where structured output consistency is paramount.
Our work has several limitations. We focus exclusively on JSON format and evaluate only six models with 75 samples, potentially missing patterns in other structured formats or newer models. The semantic similarity metrics may not capture domain-specific equivalences, and the computational cost of multiple generations for consistency evaluation may limit adoption. Future work should extend to other formats, expand model coverage, and optimize evaluation efficiency.

\bibliographystyle{unsrt}
\bibliography{sted_references}

% After your main content and references
\appendix
\section{Dataset Characteristics and Distribution Analysis}
\label{sec:dataset-distribution}

We provide a comprehensive analysis of the 75 base samples used to generate our 2,400 synthetic test cases. This analysis demonstrates the representativeness and diversity of our dataset across multiple dimensions of JSON structural complexity.

\subsection{Structural Complexity Distribution}
\label{subsec:structural-complexity}

Table~\ref{tab:json-depth-appendix} shows the distribution of JSON depths across our base samples, with the majority (58.7\%) at depth 4, aligning with typical enterprise API structures.

\begin{table}[H]
\centering
\caption{JSON Depth Distribution Across Base Samples}
\label{tab:json-depth-appendix}
\begin{tabular}{lccp{5cm}}
\toprule
\textbf{Depth Level} & \textbf{Count} & \textbf{Percentage} & \textbf{Interpretation} \\
\midrule
Depth 2 & 8 & 10.7\% & Simple, flat structures \\
Depth 3 & 7 & 9.3\% & Moderately nested \\
\textbf{Depth 4} & \textbf{44} & \textbf{58.7\%} & \textbf{Most common depth} \\
Depth 5 & 13 & 17.3\% & Complex nested structures \\
Depth 6 & 2 & 2.7\% & Highly complex \\
Depth 7 & 1 & 1.3\% & Maximum complexity \\
\bottomrule
\end{tabular}
\end{table}

Table~\ref{tab:field-count-appendix} presents the field count distribution, demonstrating coverage from simple configurations to complex documents.

\begin{table}[H]
\centering
\caption{Field Count Distribution}
\label{tab:field-count-appendix}
\begin{tabular}{lccp{4.5cm}}
\toprule
\textbf{Field Range} & \textbf{Count} & \textbf{Percentage} & \textbf{Real-world Analogy} \\
\midrule
1--10 fields & 6 & 8.0\% & Simple config files \\
11--25 fields & 22 & 29.3\% & API responses \\
\textbf{26--50 fields} & \textbf{27} & \textbf{36.0\%} & \textbf{Most common} \\
51--100 fields & 16 & 21.3\% & Complex documents \\
100+ fields & 4 & 5.3\% & Large schemas \\
\bottomrule
\end{tabular}
\end{table}

\subsection{Field Type and Complexity Analysis}
\label{subsec:field-type-analysis}

Tables~\ref{tab:field-types-appendix} and \ref{tab:complexity-summary-appendix} provide detailed breakdowns of field types and complexity metrics across the dataset.

\begin{table}[H]
\centering
\caption{Distribution of Field Types Across All Samples}
\label{tab:field-types-appendix}
\begin{tabular}{lrrp{5.5cm}}
\toprule
\textbf{Field Type} & \textbf{Count} & \textbf{Percentage} & \textbf{Coverage} \\
\midrule
String & 7,276 & 68.2\% & Text data, IDs, descriptions \\
Integer & 1,968 & 18.4\% & Counts, IDs, numeric values \\
Array & 907 & 8.5\% & Lists, collections \\
Object & 524 & 4.9\% & Nested structures \\
\bottomrule
\end{tabular}
\end{table}

\begin{table}[H]
\centering
\caption{Statistical Summary of Complexity Metrics}
\label{tab:complexity-summary-appendix}
\begin{tabular}{lrrrrl}
\toprule
\textbf{Metric} & \textbf{Min} & \textbf{Max} & \textbf{Mean} & \textbf{Std Dev} & \textbf{Interpretation} \\
\midrule
Max Depth & 2 & 7 & 4.0 & 1.0 & Good variety \\
Total Fields & 4 & 228 & 42.3 & 37.8 & Wide range \\
Total Nodes & 10 & 320 & 64.2 & 44.9 & Diverse complexity \\
Arrays & 0 & 11 & 4.7 & 2.7 & Adequate coverage \\
Nested Objects & 0 & 98 & 12.1 & 13.6 & Strong variety \\
\bottomrule
\end{tabular}
\end{table}

\section{Statistical Analysis Details}
\label{sec:statistical-analysis}

\subsection{Methodology}

We performed comprehensive statistical validation using:
\begin{itemize}
\item \textbf{Paired t-tests}: Each of the 75 base samples was evaluated by all metrics on the same variations, enabling paired comparisons
\item \textbf{Bonferroni correction}: Applied for multiple comparisons ($\alpha = 0.05/30 = 0.0017$)
\item \textbf{Effect size calculation}: Cohen's d to quantify practical significance
\item \textbf{Confidence intervals}: 95\% CIs computed using bootstrap with 1000 iterations
\end{itemize}

\subsection{Semantic Variations - Detailed Results}

\begin{table}[H]
\centering
\caption{Summary Statistics for Semantic Variations (N=750)}
\label{tab:semantic-stats-detailed}
\begin{tabular}{lcccccc}
\toprule
\textbf{Metric} & \textbf{Mean} & \textbf{Std Dev} & \textbf{95\% CI} & \textbf{Min} & \textbf{Max} & \textbf{Median} \\
\midrule
STED & 0.9539 & 0.0393 & [0.9511, 0.9567] & 0.832 & 1.000 & 0.961 \\
BERTScore & 0.9582 & 0.0249 & [0.9564, 0.9600] & 0.871 & 1.000 & 0.963 \\
TED & 1.0000 & 0.0000 & [1.0000, 1.0000] & 1.000 & 1.000 & 1.000 \\
DeepDiff & 0.7991 & 0.1647 & [0.7873, 0.8109] & 0.412 & 1.000 & 0.825 \\
\bottomrule
\end{tabular}
\end{table}

\begin{table}[H]
\centering
\caption{Pairwise Statistical Comparisons - Semantic Variations}
\label{tab:semantic-pairwise}
\small
\begin{tabular}{lccccccccccc}
\toprule
\textbf{STED vs} & \multicolumn{10}{c}{\textbf{p-values by Variation Ratio}} & \textbf{Mean p} \\
\cmidrule{2-11}
& 0.1 & 0.2 & 0.3 & 0.4 & 0.5 & 0.6 & 0.7 & 0.8 & 0.9 & 1.0 & \\
\midrule
TED & <.001 & <.001 & <.001 & <.001 & <.001 & <.001 & <.001 & <.001 & <.001 & <.001 & <.001*** \\
BERTScore & .391 & .197 & .724 & .721 & .857 & .973 & .655 & .575 & .454 & .454 & .600 (ns) \\
DeepDiff & <.001 & <.001 & <.001 & <.001 & <.001 & <.001 & <.001 & <.001 & <.001 & <.001 & <.001*** \\
\bottomrule
\end{tabular}
\vspace{0.3em}
\small{*** Significant after Bonferroni correction (p < 0.0017), ns = not significant}
\end{table}

\textbf{Interpretation}: STED and BERTScore show statistically equivalent performance on semantic variations (mean p = 0.600), both correctly maintaining high similarity for semantic equivalents. TED's complete insensitivity (constant 1.0) and DeepDiff's over-sensitivity (mean 0.799) are both significantly different from STED (p < 0.001).

\subsection{Expression Variations - Detailed Results}

\begin{table}[H]
\centering
\caption{Summary Statistics for Expression Variations (N=750)}
\label{tab:expression-stats-detailed}
\begin{tabular}{lcccccc}
\toprule
\textbf{Metric} & \textbf{Mean} & \textbf{Std Dev} & \textbf{95\% CI} & \textbf{Min} & \textbf{Max} & \textbf{Median} \\
\midrule
STED & 0.9812 & 0.0174 & [0.9799, 0.9824] & 0.918 & 1.000 & 0.985 \\
BERTScore & 0.9595 & 0.0267 & [0.9576, 0.9614] & 0.882 & 1.000 & 0.964 \\
TED & 1.0000 & 0.0000 & [1.0000, 1.0000] & 1.000 & 1.000 & 1.000 \\
DeepDiff & 0.8051 & 0.1644 & [0.7934, 0.8169] & 0.423 & 1.000 & 0.831 \\
\bottomrule
\end{tabular}
\end{table}

\begin{table}[H]
\centering
\caption{Pairwise Statistical Comparisons - Expression Variations}
\label{tab:expression-pairwise}
\small
\begin{tabular}{lccccccccccc}
\toprule
\textbf{STED vs} & \multicolumn{10}{c}{\textbf{p-values by Variation Ratio}} & \textbf{Mean p} \\
\cmidrule{2-11}
& 0.1 & 0.2 & 0.3 & 0.4 & 0.5 & 0.6 & 0.7 & 0.8 & 0.9 & 1.0 & \\
\midrule
TED & <.001 & <.001 & <.001 & <.001 & <.001 & <.001 & <.001 & <.001 & <.001 & <.001 & <.001*** \\
BERTScore & <.001 & <.001 & <.001 & <.001 & <.001 & <.001 & <.001 & <.001 & <.001 & <.001 & <.001*** \\
DeepDiff & <.001 & <.001 & <.001 & <.001 & <.001 & <.001 & <.001 & <.001 & <.001 & <.001 & <.001*** \\
\bottomrule
\end{tabular}
\vspace{0.3em}
\small{*** Significant after Bonferroni correction (p < 0.0017)}
\end{table}

\textbf{Interpretation}: STED significantly outperforms all baselines on expression variations (p < 0.001), demonstrating superior format-agnostic understanding. The lower standard deviation (0.0174) compared to semantic variations (0.0393) indicates more consistent performance.

\subsection{Effect Size Analysis}

\begin{table}[H]
\centering
\caption{Cohen's d Effect Sizes for STED Comparisons}
\label{tab:effect-sizes}
\begin{tabular}{lcccc}
\toprule
\textbf{Comparison} & \multicolumn{2}{c}{\textbf{Semantic Variations}} & \multicolumn{2}{c}{\textbf{Expression Variations}} \\
\cmidrule{2-3} \cmidrule{4-5}
& Cohen's d & Interpretation & Cohen's d & Interpretation \\
\midrule
STED vs TED & -1.23 & Large & -1.45 & Large \\
STED vs BERTScore & -0.12 & Negligible & 0.82 & Large \\
STED vs DeepDiff & 1.26 & Large & 1.38 & Large \\
\bottomrule
\end{tabular}
\vspace{0.3em}
\small{Cohen's d interpretation: 0.2 = small, 0.5 = medium, 0.8 = large}
\end{table}

\subsection{Statistical Power Analysis}

With 75 samples per variation ratio and 10 ratios per variation type:
\begin{itemize}
\item \textbf{Statistical power}: 0.99 for detecting medium effect sizes (d = 0.5) at $\alpha = 0.05$
\item \textbf{Minimum detectable effect}: d = 0.33 with 80\% power
\item \textbf{Sample size adequacy}: Our 750 samples per variation type exceed requirements for robust conclusions
\end{itemize}

\subsection{Schema Variation Robustness}

We evaluate metric robustness under common schema evolution scenarios that preserve semantic equivalence but alter structural representation.

\subsubsection{Field Name Evolution}

Real-world schemas evolve through refactoring, where field names change while preserving meaning (e.g., API versioning). We simulate this by progressively renaming fields to semantically equivalent alternatives.

\begin{table}[H]
\centering
\caption{Metric robustness to schema field renaming (N=75)}
\label{tab:field-name-robustness}
\begin{tabular}{lcccc}
\toprule
\textbf{Fields} & \textbf{TED} & \textbf{STED} & \textbf{BERTScore} & \textbf{DeepDiff} \\
\textbf{Renamed} & & & & \\
\midrule
10\% & 0.942 & 0.903 & 0.996 & 0.947 \\
20\% & 0.915 & 0.893 & 0.995 & 0.942 \\
30\% & 0.900 & 0.886 & 0.994 & 0.938 \\
40\% & 0.879 & 0.882 & 0.993 & 0.933 \\
50\% & 0.848 & 0.877 & 0.992 & 0.928 \\
60\% & 0.831 & 0.874 & 0.991 & 0.923 \\
70\% & 0.811 & 0.870 & 0.990 & 0.919 \\
80\% & 0.790 & 0.866 & 0.989 & 0.916 \\
90\% & 0.775 & 0.862 & 0.989 & 0.915 \\
100\% & 0.759 & 0.856 & 0.988 & 0.913 \\
\midrule
\textbf{Degradation}$^a$ & -19.4\% & -5.2\% & -0.8\% & -3.6\% \\
\bottomrule
\end{tabular}
\vspace{-2mm}
\begin{flushleft}
\small $^a$ Relative decrease from 0\% renamed (baseline=1.0) to 100\% renamed
\end{flushleft}
\end{table}

\subsubsection{Schema Restructuring Patterns}

We test two common schema refactoring patterns that preserve information content:

\begin{table}[H]
\centering
\caption{Metric sensitivity to schema restructuring patterns}
\label{tab:structural-perturbations}
\begin{tabular}{lcc}
\toprule
\textbf{Metric} & \textbf{Schema Flattening}$^a$ & \textbf{Schema Nesting}$^b$ \\
\midrule
TED & 0.000 & 0.506 \\
STED & 0.051 & 0.000 \\
BERTScore & 0.897 & 0.989 \\
DeepDiff & 0.809 & 0.781 \\
\bottomrule
\end{tabular}
\vspace{2mm}
\begin{flushleft}
\small $^a$ \textbf{Flattening}: Denormalizing nested objects into flat structure\\
\small Example: \texttt{\{"user": \{"name": "John", "age": 30\}\}} → \texttt{\{"user\_name": "John", "user\_age": 30\}}\\[1mm]
\small $^b$ \textbf{Nesting}: Organizing flat fields into logical groups\\
\small Example: \texttt{\{"street": "Main", "city": "NYC"\}} → \texttt{\{"address": \{"street": "Main", "city": "NYC"\}\}}
\end{flushleft}
\end{table}

The results demonstrate that:
\begin{itemize}
\item \textbf{TED} shows steep degradation (-19.4\%) as field names change, treating renamed fields as completely different
\item \textbf{STED} maintains better robustness (-5.2\% degradation) through structural pattern recognition
\item \textbf{BERTScore} demonstrates exceptional robustness (-0.8\%) via semantic understanding of field names
\item \textbf{DeepDiff} shows moderate robustness (-3.6\%) with consistent degradation pattern
\end{itemize}

For structural transformations, TED and STED show complementary sensitivities: TED detects flattening as complete change (0.000) while STED identifies nesting as significant (0.000). This reflects their different approaches to structural similarity.

\section{Detailed Consistency Results}

\subsection{Temperature Sensitivity Analysis}

We analyze how temperature settings affect model consistency across 9 temperature values from 0.1 to 0.9. Each cell represents mean consistency score across 80 test cases (720 total per model).

\begin{table}[H]
\centering
\caption{Overall consistency scores across temperature settings}
\label{tab:overall-consistency-temp}
\begin{tabular}{lrrrrrrrrr}
\toprule
\textbf{Model} & \multicolumn{9}{c}{\textbf{Temperature}} \\
\cmidrule{2-10}
 & 0.1 & 0.2 & 0.3 & 0.4 & 0.5 & 0.6 & 0.7 & 0.8 & 0.9 \\
\midrule
Claude-3.5-Sonnet-V2 & 0.595 & 0.666 & 0.628 & 0.607 & 0.603 & 0.604 & 0.577 & 0.578 & 0.570 \\
Claude-3-Haiku & 0.527 & 0.479 & 0.490 & 0.433 & 0.432 & 0.344 & 0.371 & 0.374 & 0.361 \\
Claude-3.5-Haiku & 0.674 & 0.650 & 0.537 & 0.530 & 0.513 & 0.497 & 0.448 & 0.442 & 0.428 \\
Claude-3.7-Sonnet & \textbf{0.717} & \textbf{0.698} & \textbf{0.690} & \textbf{0.672} & \textbf{0.661} & \textbf{0.648} & \textbf{0.623} & \textbf{0.628} & \textbf{0.581} \\
Llama-3.3-70B & 0.519 & 0.522 & 0.457 & 0.411 & 0.346 & 0.386 & 0.368 & 0.374 & 0.334 \\
Nova-Pro & 0.438 & 0.507 & 0.454 & 0.457 & 0.494 & 0.448 & 0.380 & 0.420 & 0.401 \\
\bottomrule
\end{tabular}
\end{table}

\begin{table}[H]
\centering
\caption{Semantic consistency scores across temperature settings}
\label{tab:semantic-consistency-temp}
\begin{tabular}{lrrrrrrrrr}
\toprule
\textbf{Model} & \multicolumn{9}{c}{\textbf{Temperature}} \\
\cmidrule{2-10}
 & 0.1 & 0.2 & 0.3 & 0.4 & 0.5 & 0.6 & 0.7 & 0.8 & 0.9 \\
\midrule
Claude-3.5-Sonnet-V2 & 0.464 & 0.535 & 0.481 & 0.464 & 0.463 & 0.450 & 0.423 & 0.432 & 0.416 \\
Claude-3-Haiku & 0.456 & 0.406 & 0.402 & 0.337 & 0.336 & 0.263 & 0.283 & 0.279 & 0.264 \\
Claude-3.5-Haiku & \textbf{0.583} & 0.537 & 0.443 & 0.424 & 0.427 & 0.380 & 0.354 & 0.339 & 0.319 \\
Claude-3.7-Sonnet & 0.565 & \textbf{0.547} & \textbf{0.535} & \textbf{0.508} & \textbf{0.498} & \textbf{0.476} & \textbf{0.454} & \textbf{0.462} & \textbf{0.414} \\
Llama-3.3-70B & 0.387 & 0.390 & 0.326 & 0.292 & 0.222 & 0.263 & 0.246 & 0.257 & 0.215 \\
Nova-Pro & 0.342 & 0.404 & 0.354 & 0.348 & 0.377 & 0.338 & 0.283 & 0.314 & 0.300 \\
\bottomrule
\end{tabular}
\end{table}

\begin{table}[H]
\centering
\caption{Structural consistency scores across temperature settings}
\label{tab:structural-consistency-temp}
\begin{tabular}{lrrrrrrrrr}
\toprule
\textbf{Model} & \multicolumn{9}{c}{\textbf{Temperature}} \\
\cmidrule{2-10}
 & 0.1 & 0.2 & 0.3 & 0.4 & 0.5 & 0.6 & 0.7 & 0.8 & 0.9 \\
\midrule
Claude-3.5-Sonnet-V2 & 0.875 & 0.975 & 0.963 & 0.950 & 0.950 & 0.963 & 0.950 & 0.938 & 0.950 \\
Claude-3-Haiku & 0.755 & 0.750 & 0.781 & 0.767 & 0.758 & 0.679 & 0.682 & 0.708 & 0.722 \\
Claude-3.5-Haiku & 0.860 & 0.869 & 0.776 & 0.768 & 0.757 & 0.771 & 0.703 & 0.722 & 0.704 \\
Claude-3.7-Sonnet & \textbf{1.000} & \textbf{1.000} & \textbf{1.000} & \textbf{1.000} & \textbf{1.000} & \textbf{1.000} & \textbf{1.000} & \textbf{1.000} & \textbf{0.988} \\
Llama-3.3-70B & 0.825 & 0.863 & 0.825 & 0.800 & 0.800 & 0.850 & 0.825 & 0.838 & 0.838 \\
Nova-Pro & 0.722 & 0.775 & 0.700 & 0.739 & 0.787 & 0.749 & 0.658 & 0.748 & 0.673 \\
\bottomrule
\end{tabular}
\end{table}

\subsection{Summary Statistics}

\begin{table}[H]
\centering
\caption{Aggregated consistency metrics (Mean ± SD across all temperatures, N=720 per model)}
\label{tab:consistency-summary}
\begin{tabular}{lccc}
\toprule
\textbf{Model} & \textbf{Overall} & \textbf{Semantic} & \textbf{Structural} \\
\midrule
Claude-3.5-Sonnet-V2 & 0.603±0.303 & 0.459±0.318 & 0.946±0.227 \\
Claude-3-Haiku & 0.424±0.386 & 0.336±0.399 & 0.734±0.432 \\
Claude-3.5-Haiku & 0.524±0.368 & 0.423±0.378 & 0.770±0.411 \\
Claude-3.7-Sonnet & \textbf{0.658±0.240} & \textbf{0.495±0.288} & \textbf{0.999±0.037} \\
Llama-3.3-70B & 0.413±0.337 & 0.289±0.340 & 0.829±0.377 \\
Nova-Pro & 0.444±0.378 & 0.340±0.388 & 0.728±0.430 \\
\bottomrule
\end{tabular}
\end{table}

\subsection{Key Findings}

\paragraph{Temperature Effects}
\begin{itemize}
\item Most models show declining consistency as temperature increases, with steepest drops between T=0.5 and T=0.7
\item Claude-3.5-Haiku exhibits strongest temperature sensitivity (46\% decline from T=0.1 to T=0.9)
\item Claude-3.7-Sonnet maintains most stable performance across temperatures (19\% decline)
\item Structural consistency remains more robust to temperature changes than semantic consistency
\end{itemize}

\paragraph{Optimal Temperature Ranges}
\begin{itemize}
\item \textbf{T=0.1-0.3}: Best for consistency, all models achieve peak performance
\item \textbf{T=0.4-0.6}: Moderate degradation, acceptable for most applications
\item \textbf{T=0.7-0.9}: Significant consistency loss, especially for semantic preservation
\end{itemize}

\paragraph{Model-Specific Observations}
\begin{itemize}
\item \textbf{Claude-3.7-Sonnet}: Near-perfect structural consistency (≥0.988) across all temperatures
\item \textbf{Nova-Pro}: Shows irregular pattern with local maximum at T=0.5
\item \textbf{Llama-3.3-70B}: Exhibits sharp semantic consistency drop at T=0.5
\end{itemize}

\newpage
\section*{NeurIPS Paper Checklist}

\begin{enumerate}

\item {\bf Claims}
    \item[] Question: Do the main claims made in the abstract and introduction accurately reflect the paper's contributions and scope?
    \item[] Answer: \answerYes{} % Replace by \answerYes{}, \answerNo{}, or \answerNA{}.
    \item[] Justification: \justification The abstract and introduction clearly state our three main contributions: 1.a comprehensive evaluation framework for assessing JSON consistency in LLMs across structural and semantic dimensions; 2.empirical analysis of six state-of-the-art models revealing distinct behavioral patterns, and 3. practical insights for model selection based on task requirements{}
    \item[] Guidelines:
    \begin{itemize}
        \item The answer NA means that the abstract and introduction do not include the claims made in the paper.
        \item The abstract and/or introduction should clearly state the claims made, including the contributions made in the paper and important assumptions and limitations. A No or NA answer to this question will not be perceived well by the reviewers. 
        \item The claims made should match theoretical and experimental results, and reflect how much the results can be expected to generalize to other settings. 
        \item It is fine to include aspirational goals as motivation as long as it is clear that these goals are not attained by the paper. 
    \end{itemize}

\item {\bf Limitations}
    \item[] Question: Does the paper discuss the limitations of the work performed by the authors?
    \item[] Answer: \answerYes{} % Replace by \answerYes{}, \answerNo{}, or \answerNA{}.
    \item[] Justification: \justification Discussed the limitation in conclusion section{}
    \item[] Guidelines:
    \begin{itemize}
        \item The answer NA means that the paper has no limitation while the answer No means that the paper has limitations, but those are not discussed in the paper. 
        \item The authors are encouraged to create a separate "Limitations" section in their paper.
        \item The paper should point out any strong assumptions and how robust the results are to violations of these assumptions (e.g., independence assumptions, noiseless settings, model well-specification, asymptotic approximations only holding locally). The authors should reflect on how these assumptions might be violated in practice and what the implications would be.
        \item The authors should reflect on the scope of the claims made, e.g., if the approach was only tested on a few datasets or with a few runs. In general, empirical results often depend on implicit assumptions, which should be articulated.
        \item The authors should reflect on the factors that influence the performance of the approach. For example, a facial recognition algorithm may perform poorly when image resolution is low or images are taken in low lighting. Or a speech-to-text system might not be used reliably to provide closed captions for online lectures because it fails to handle technical jargon.
        \item The authors should discuss the computational efficiency of the proposed algorithms and how they scale with dataset size.
        \item If applicable, the authors should discuss possible limitations of their approach to address problems of privacy and fairness.
        \item While the authors might fear that complete honesty about limitations might be used by reviewers as grounds for rejection, a worse outcome might be that reviewers discover limitations that aren't acknowledged in the paper. The authors should use their best judgment and recognize that individual actions in favor of transparency play an important role in developing norms that preserve the integrity of the community. Reviewers will be specifically instructed to not penalize honesty concerning limitations.
    \end{itemize}

\item {\bf Theory assumptions and proofs}
    \item[] Question: For each theoretical result, does the paper provide the full set of assumptions and a complete (and correct) proof?
    \item[] Answer: \answerNA{} % Replace by \answerYes{}, \answerNo{}, or \answerNA{}.
    \item[] Justification: \justification This paper presents an empirical evaluation framework and metric design for LLM-generated structured outputs. It does not include theoretical results, theorems, lemmas, or formal proofs. The work focuses on experimental validation of the proposed STED metric and consistency score through empirical benchmarking and comparative analysis with existing methods{}
    \item[] Guidelines:
    \begin{itemize}
        \item The answer NA means that the paper does not include theoretical results. 
        \item All the theorems, formulas, and proofs in the paper should be numbered and cross-referenced.
        \item All assumptions should be clearly stated or referenced in the statement of any theorems.
        \item The proofs can either appear in the main paper or the supplemental material, but if they appear in the supplemental material, the authors are encouraged to provide a short proof sketch to provide intuition. 
        \item Inversely, any informal proof provided in the core of the paper should be complemented by formal proofs provided in appendix or supplemental material.
        \item Theorems and Lemmas that the proof relies upon should be properly referenced. 
    \end{itemize}

    \item {\bf Experimental result reproducibility}
    \item[] Question: Does the paper fully disclose all the information needed to reproduce the main experimental results of the paper to the extent that it affects the main claims and/or conclusions of the paper (regardless of whether the code and data are provided or not)?
    \item[] Answer: \answerYes{} % Replace by \answerYes{}, \answerNo{}, or \answerNA{}.
    \item[] Justification: \justification The paper provides complete specifications for reproduction: STED algorithm formulation (Section 3, Equations 1-4), experimental setup with 80 synthetic samples, six named models from Amazon Bedrock, temperature ranges (0.1-0.9), 10 generations per configuration, and all baseline metrics. All models are publicly accessible through Amazon Bedrock API, and the synthetic data generation process with controlled modifications is fully described.{}
    \item[] Guidelines:
    \begin{itemize}
        \item The answer NA means that the paper does not include experiments.
        \item If the paper includes experiments, a No answer to this question will not be perceived well by the reviewers: Making the paper reproducible is important, regardless of whether the code and data are provided or not.
        \item If the contribution is a dataset and/or model, the authors should describe the steps taken to make their results reproducible or verifiable. 
        \item Depending on the contribution, reproducibility can be accomplished in various ways. For example, if the contribution is a novel architecture, describing the architecture fully might suffice, or if the contribution is a specific model and empirical evaluation, it may be necessary to either make it possible for others to replicate the model with the same dataset, or provide access to the model. In general. releasing code and data is often one good way to accomplish this, but reproducibility can also be provided via detailed instructions for how to replicate the results, access to a hosted model (e.g., in the case of a large language model), releasing of a model checkpoint, or other means that are appropriate to the research performed.
        \item While NeurIPS does not require releasing code, the conference does require all submissions to provide some reasonable avenue for reproducibility, which may depend on the nature of the contribution. For example
        \begin{enumerate}
            \item If the contribution is primarily a new algorithm, the paper should make it clear how to reproduce that algorithm.
            \item If the contribution is primarily a new model architecture, the paper should describe the architecture clearly and fully.
            \item If the contribution is a new model (e.g., a large language model), then there should either be a way to access this model for reproducing the results or a way to reproduce the model (e.g., with an open-source dataset or instructions for how to construct the dataset).
            \item We recognize that reproducibility may be tricky in some cases, in which case authors are welcome to describe the particular way they provide for reproducibility. In the case of closed-source models, it may be that access to the model is limited in some way (e.g., to registered users), but it should be possible for other researchers to have some path to reproducing or verifying the results.
        \end{enumerate}
    \end{itemize}


\item {\bf Open access to data and code}
    \item[] Question: Does the paper provide open access to the data and code, with sufficient instructions to faithfully reproduce the main experimental results, as described in supplemental material?
    \item[] Answer: \answerNo{} % Replace by \answerYes{}, \answerNo{}, or \answerNA{}.
    \item[] Justification: \justification While the paper provides detailed algorithmic descriptions and experimental setup information sufficient for reproduction, we do not currently provide open access to the implementation code or synthetic dataset. The STED algorithm is fully specified in Section 3 with mathematical formulations, and the experimental protocol is detailed in Section 4. The models used are publicly accessible through Amazon Bedrock API. We plan to release the code implementation and synthetic dataset upon paper acceptance to maintain anonymity during review.{}
    \item[] Guidelines:
    \begin{itemize}
        \item The answer NA means that paper does not include experiments requiring code.
        \item Please see the NeurIPS code and data submission guidelines (\url{https://nips.cc/public/guides/CodeSubmissionPolicy}) for more details.
        \item While we encourage the release of code and data, we understand that this might not be possible, so “No” is an acceptable answer. Papers cannot be rejected simply for not including code, unless this is central to the contribution (e.g., for a new open-source benchmark).
        \item The instructions should contain the exact command and environment needed to run to reproduce the results. See the NeurIPS code and data submission guidelines (\url{https://nips.cc/public/guides/CodeSubmissionPolicy}) for more details.
        \item The authors should provide instructions on data access and preparation, including how to access the raw data, preprocessed data, intermediate data, and generated data, etc.
        \item The authors should provide scripts to reproduce all experimental results for the new proposed method and baselines. If only a subset of experiments are reproducible, they should state which ones are omitted from the script and why.
        \item At submission time, to preserve anonymity, the authors should release anonymized versions (if applicable).
        \item Providing as much information as possible in supplemental material (appended to the paper) is recommended, but including URLs to data and code is permitted.
    \end{itemize}


\item {\bf Experimental setting/details}
    \item[] Question: Does the paper specify all the training and test details (e.g., data splits, hyperparameters, how they were chosen, type of optimizer, etc.) necessary to understand the results?
    \item[] Answer: \answerYes{} % Replace by \answerYes{}, \answerNo{}, or \answerNA{}.
    \item[] Justification: \justification Section 4 provides comprehensive experimental details. We specify the dataset (80 synthetic JSON samples with controlled modifications), all model configurations (six models from Amazon Bedrock: Claude-3.5-Sonnet-V2, Claude-3.7-Sonnet, Claude-3.5-Haiku, Claude-3-Haiku, Nova-Pro, Nova-Lite), temperature settings (0.1 to 0.9 in 0.1 increments), generation protocol (10 outputs per temperature setting), and evaluation metrics (STED with α=0.7, TED, BERTScore with deberta-large-mnli, DeepDiff). The paper clearly describes the three modification types (structural, semantic, combined) and their implementation. Since our work evaluates pre-trained models without training, no training details or optimizers are applicable.{}
    \item[] Guidelines:
    \begin{itemize}
        \item The answer NA means that the paper does not include experiments.
        \item The experimental setting should be presented in the core of the paper to a level of detail that is necessary to appreciate the results and make sense of them.
        \item The full details can be provided either with the code, in appendix, or as supplemental material.
    \end{itemize}

\item {\bf Experiment statistical significance}
    \item[] Question: Does the paper report error bars suitably and correctly defined or other appropriate information about the statistical significance of the experiments?
    \item[] Answer: \answerNo{} % Replace by \answerYes{}, \answerNo{}, or \answerNA{}.
    \item[] Justification: \justification Our experiments generate 10 outputs per temperature setting to compute consistency scores, but we report aggregate metrics without error bars or confidence intervals. Figure 2 and Table 2 present mean consistency values without variance measures or statistical significance tests for model comparisons{}
    \item[] Guidelines:
    \begin{itemize}
        \item The answer NA means that the paper does not include experiments.
        \item The authors should answer "Yes" if the results are accompanied by error bars, confidence intervals, or statistical significance tests, at least for the experiments that support the main claims of the paper.
        \item The factors of variability that the error bars are capturing should be clearly stated (for example, train/test split, initialization, random drawing of some parameter, or overall run with given experimental conditions).
        \item The method for calculating the error bars should be explained (closed form formula, call to a library function, bootstrap, etc.)
        \item The assumptions made should be given (e.g., Normally distributed errors).
        \item It should be clear whether the error bar is the standard deviation or the standard error of the mean.
        \item It is OK to report 1-sigma error bars, but one should state it. The authors should preferably report a 2-sigma error bar than state that they have a 96\% CI, if the hypothesis of Normality of errors is not verified.
        \item For asymmetric distributions, the authors should be careful not to show in tables or figures symmetric error bars that would yield results that are out of range (e.g. negative error rates).
        \item If error bars are reported in tables or plots, The authors should explain in the text how they were calculated and reference the corresponding figures or tables in the text.
    \end{itemize}

\item {\bf Experiments compute resources}
    \item[] Question: For each experiment, does the paper provide sufficient information on the computer resources (type of compute workers, memory, time of execution) needed to reproduce the experiments?
    \item[] Answer: \answerNo{} % Replace by \answerYes{}, \answerNo{}, or \answerNA{}.
    \item[] Justification: \justification While we specify that models are accessed through Amazon Bedrock API, we do not provide details on compute resources, execution time, or memory requirements. The paper does not report the time needed to generate 10 outputs per temperature setting across six models, nor the computational cost of STED calculations on 80 samples.{}
    \item[] Guidelines:
    \begin{itemize}
        \item The answer NA means that the paper does not include experiments.
        \item The paper should indicate the type of compute workers CPU or GPU, internal cluster, or cloud provider, including relevant memory and storage.
        \item The paper should provide the amount of compute required for each of the individual experimental runs as well as estimate the total compute. 
        \item The paper should disclose whether the full research project required more compute than the experiments reported in the paper (e.g., preliminary or failed experiments that didn't make it into the paper). 
    \end{itemize}
    
\item {\bf Code of ethics}
    \item[] Question: Does the research conducted in the paper conform, in every respect, with the NeurIPS Code of Ethics \url{https://neurips.cc/public/EthicsGuidelines}?
    \item[] Answer: \answerNo{} % Replace by \answerYes{}, \answerNo{}, or \answerNA{}.
    \item[] Justification: \justification Our research on evaluating LLM consistency for structured outputs adheres to all aspects of the NeurIPS Code of Ethics. The work involves no human subjects, uses publicly available models through commercial APIs (Amazon Bedrock), generates only synthetic data for evaluation purposes, and presents results objectively without misleading claims. The research aims to improve reliability and transparency of LLM systems, contributing positively to the field.{}
    \item[] Guidelines:
    \begin{itemize}
        \item The answer NA means that the authors have not reviewed the NeurIPS Code of Ethics.
        \item If the authors answer No, they should explain the special circumstances that require a deviation from the Code of Ethics.
        \item The authors should make sure to preserve anonymity (e.g., if there is a special consideration due to laws or regulations in their jurisdiction).
    \end{itemize}


\item {\bf Broader impacts}
    \item[] Question: Does the paper discuss both potential positive societal impacts and negative societal impacts of the work performed?
    \item[] Answer: \answerNo{} % Replace by \answerYes{}, \answerNo{}, or \answerNA{}.
    \item[] Justification: \justification Our work on consistency metrics for LLM-generated structured outputs is foundational research focused on evaluation methodology. The primary impact is improving reliability of LLM systems in production environments, which has clear positive implications for system dependability. While more consistent LLMs could theoretically be misused for generating convincing disinformation at scale, this risk exists for any LLM improvement and is not specific to consistency evaluation. The metric itself is a diagnostic tool without direct negative applications.{}
    \item[] Guidelines:
    \begin{itemize}
        \item The answer NA means that there is no societal impact of the work performed.
        \item If the authors answer NA or No, they should explain why their work has no societal impact or why the paper does not address societal impact.
        \item Examples of negative societal impacts include potential malicious or unintended uses (e.g., disinformation, generating fake profiles, surveillance), fairness considerations (e.g., deployment of technologies that could make decisions that unfairly impact specific groups), privacy considerations, and security considerations.
        \item The conference expects that many papers will be foundational research and not tied to particular applications, let alone deployments. However, if there is a direct path to any negative applications, the authors should point it out. For example, it is legitimate to point out that an improvement in the quality of generative models could be used to generate deepfakes for disinformation. On the other hand, it is not needed to point out that a generic algorithm for optimizing neural networks could enable people to train models that generate Deepfakes faster.
        \item The authors should consider possible harms that could arise when the technology is being used as intended and functioning correctly, harms that could arise when the technology is being used as intended but gives incorrect results, and harms following from (intentional or unintentional) misuse of the technology.
        \item If there are negative societal impacts, the authors could also discuss possible mitigation strategies (e.g., gated release of models, providing defenses in addition to attacks, mechanisms for monitoring misuse, mechanisms to monitor how a system learns from feedback over time, improving the efficiency and accessibility of ML).
    \end{itemize}
    
\item {\bf Safeguards}
    \item[] Question: Does the paper describe safeguards that have been put in place for responsible release of data or models that have a high risk for misuse (e.g., pretrained language models, image generators, or scraped datasets)?
    \item[] Answer: \answerNA{} % Replace by \answerYes{}, \answerNo{}, or \answerNA{}.
    \item[] Justification: \justification Our paper presents an evaluation metric (STED) and synthetic dataset for testing LLM consistency. We do not release any pretrained models, image generators, or scraped datasets. The synthetic JSON data used for evaluation contains no personal information or potentially harmful content. The STED metric is a measurement tool that poses no risk for misuse. We evaluate existing commercial models accessed through Amazon Bedrock's API, which already has its own safeguards and usage policies.{}
    \item[] Guidelines:
    \begin{itemize}
        \item The answer NA means that the paper poses no such risks.
        \item Released models that have a high risk for misuse or dual-use should be released with necessary safeguards to allow for controlled use of the model, for example by requiring that users adhere to usage guidelines or restrictions to access the model or implementing safety filters. 
        \item Datasets that have been scraped from the Internet could pose safety risks. The authors should describe how they avoided releasing unsafe images.
        \item We recognize that providing effective safeguards is challenging, and many papers do not require this, but we encourage authors to take this into account and make a best faith effort.
    \end{itemize}

\item {\bf Licenses for existing assets}
    \item[] Question: Are the creators or original owners of assets (e.g., code, data, models), used in the paper, properly credited and are the license and terms of use explicitly mentioned and properly respected?
    \item[] Answer: \answerYes{} % Replace by \answerYes{}, \answerNo{}, or \answerNA{}.
    \item[] Justification: \justification We properly credit all existing assets used in our work. The LLMs (Claude-3.5-Sonnet-V2, Claude-3.7-Sonnet, Claude-3.5-Haiku, Claude-3-Haiku, Nova-Pro, Nova-Lite) are accessed through Amazon Bedrock API under commercial terms of service. We cite the original papers for baseline metrics: Tree Edit Distance (Zhang & Shasha, 1989), BERTScore (Zhang et al., 2020), and reference DeepDiff as an open-source Python library. The DeBERTa model used for embeddings is properly attributed. All usage complies with respective terms of service and licensing agreements.{}
    \item[] Guidelines:
    \begin{itemize}
        \item The answer NA means that the paper does not use existing assets.
        \item The authors should cite the original paper that produced the code package or dataset.
        \item The authors should state which version of the asset is used and, if possible, include a URL.
        \item The name of the license (e.g., CC-BY 4.0) should be included for each asset.
        \item For scraped data from a particular source (e.g., website), the copyright and terms of service of that source should be provided.
        \item If assets are released, the license, copyright information, and terms of use in the package should be provided. For popular datasets, \url{paperswithcode.com/datasets} has curated licenses for some datasets. Their licensing guide can help determine the license of a dataset.
        \item For existing datasets that are re-packaged, both the original license and the license of the derived asset (if it has changed) should be provided.
        \item If this information is not available online, the authors are encouraged to reach out to the asset's creators.
    \end{itemize}

\item {\bf New assets}
    \item[] Question: Are new assets introduced in the paper well documented and is the documentation provided alongside the assets?
    \item[] Answer: \answerNA{} % Replace by \answerYes{}, \answerNo{}, or \answerNA{}.
    \item[] Justification: \justification The paper does not release new assets at submission time. While we introduce the STED metric and create synthetic evaluation data, these are not currently provided as downloadable assets. The STED algorithm is fully specified in the paper (Section 3) with mathematical formulations rather than as a released codebase. The synthetic dataset is described but not released. We plan to release both the implementation and dataset upon paper acceptance.{}
    \item[] Guidelines:
    \begin{itemize}
        \item The answer NA means that the paper does not release new assets.
        \item Researchers should communicate the details of the dataset/code/model as part of their submissions via structured templates. This includes details about training, license, limitations, etc. 
        \item The paper should discuss whether and how consent was obtained from people whose asset is used.
        \item At submission time, remember to anonymize your assets (if applicable). You can either create an anonymized URL or include an anonymized zip file.
    \end{itemize}

\item {\bf Crowdsourcing and research with human subjects}
    \item[] Question: For crowdsourcing experiments and research with human subjects, does the paper include the full text of instructions given to participants and screenshots, if applicable, as well as details about compensation (if any)? 
    \item[] Answer: \answerNA{} % Replace by \answerYes{}, \answerNo{}, or \answerNA{}.
    \item[] Justification: \justification Our paper does not involve crowdsourcing or research with human subjects. All experiments use synthetic data generated programmatically and evaluate automated LLM outputs without human annotation or evaluation. The consistency metrics are computed automatically without human involvement.{}
    \item[] Guidelines:
    \begin{itemize}
        \item The answer NA means that the paper does not involve crowdsourcing nor research with human subjects.
        \item Including this information in the supplemental material is fine, but if the main contribution of the paper involves human subjects, then as much detail as possible should be included in the main paper. 
        \item According to the NeurIPS Code of Ethics, workers involved in data collection, curation, or other labor should be paid at least the minimum wage in the country of the data collector. 
    \end{itemize}

\item {\bf Institutional review board (IRB) approvals or equivalent for research with human subjects}
    \item[] Question: Does the paper describe potential risks incurred by study participants, whether such risks were disclosed to the subjects, and whether Institutional Review Board (IRB) approvals (or an equivalent approval/review based on the requirements of your country or institution) were obtained?
    \item[] Answer: \answerNA{} % Replace by \answerYes{}, \answerNo{}, or \answerNA{}.
    \item[] Justification: \justification Our paper does not involve research with human subjects or crowdsourcing. All experiments use synthetic data and automated evaluation of LLM outputs without any human participation, annotation, or assessment. Therefore, IRB approval is not applicable.{}
    \item[] Guidelines:
    \begin{itemize}
        \item The answer NA means that the paper does not involve crowdsourcing nor research with human subjects.
        \item Depending on the country in which research is conducted, IRB approval (or equivalent) may be required for any human subjects research. If you obtained IRB approval, you should clearly state this in the paper. 
        \item We recognize that the procedures for this may vary significantly between institutions and locations, and we expect authors to adhere to the NeurIPS Code of Ethics and the guidelines for their institution. 
        \item For initial submissions, do not include any information that would break anonymity (if applicable), such as the institution conducting the review.
    \end{itemize}

\item {\bf Declaration of LLM usage}
    \item[] Question: Does the paper describe the usage of LLMs if it is an important, original, or non-standard component of the core methods in this research? Note that if the LLM is used only for writing, editing, or formatting purposes and does not impact the core methodology, scientific rigorousness, or originality of the research, declaration is not required.
    %this research? 
    \item[] Answer: \answerYes{} % Replace by \answerYes{}, \answerNo{}, or \answerNA{}.
    \item[] Justification: \justification Our paper explicitly describes LLM usage as test subjects for evaluation, not as part of the methodology development. Section 4 clearly identifies six LLMs (Claude-3.5-Sonnet-V2, Claude-3.7-Sonnet, Claude-3.5-Haiku, Claude-3-Haiku, Nova-Pro, Nova-Lite) accessed via Amazon Bedrock API. These models are the evaluation targets for our consistency metrics, with their outputs being analyzed rather than contributing to the method design. The STED metric itself is developed independently of LLM assistance.{}
    \item[] Guidelines:
    \begin{itemize}
        \item The answer NA means that the core method development in this research does not involve LLMs as any important, original, or non-standard components.
        \item Please refer to our LLM policy (\url{https://neurips.cc/Conferences/2025/LLM}) for what should or should not be described.
    \end{itemize}

\end{enumerate}

\end{document}